% ============================================================================
%   CHAPTER 4 — METHODOLOGY AND TECHNICAL APPROACH
%   Aivancity: 15-20 pages — the technical core of the thesis
%   Should be reproducible and self-contained
% ============================================================================

\chapter{Methodology and Technical Approach}
\label{chap:methodology}

% ----------------------------------------------------------------------------
\section{Overview of the Experimental Design}
\label{sec:met-overview}

This chapter describes the methodology used to audit the calibration of large 
language models in patient--clinical trial matching. The experimental design 
follows a pre-registered protocol consistent with the TRIPOD+AI reporting 
guidelines~\cite{collins2024tripod}, which recommend explicit hypothesis 
specification, statistical rigor, and reproducibility for clinical prediction 
models.

The evaluation pipeline is organized around three principles. First, all 
predictions are generated on a single held-out benchmark, 
TrialGPT-Criterion-Annotations~\cite{jin2024trialgpt}, which contains 1{,}015 
patient--criterion pairs annotated by three physicians and by GPT-4. Second, 
model probabilities are extracted directly from the log-probabilities of 
candidate labels rather than from free-form generation, in order to obtain 
well-defined probability distributions suitable for calibration analysis. 
Third, every hypothesis test is defined before the analysis begins, and 
p-values are corrected for multiple comparisons.

The design forms a $2 \times 2$ matrix crossing model type and post-training 
regime. Along one axis, we compare generalist models (Mistral-7B, 
Llama-3.1-8B) with medically specialized models (Meditron-7B, BioMistral-7B). 
Along the other axis, we compare base models (without RLHF) with 
instruction-tuned models (with RLHF). This crossing produces six models, each 
of which is evaluated with two prompt variants: a generic minimalist prompt 
(Prompt~A) and an adapted prompt with expert guidance (Prompt~B). The full 
factorial yields twelve configurations, all evaluated on the same 1{,}015 
pairs, for a total of 12{,}180 predictions.

Figure~\ref{fig:met-overview} summarizes the pipeline. Section~\ref{sec:met-hypotheses} 
states the four pre-registered hypotheses. Section~\ref{sec:met-data} describes 
the dataset. Sections~\ref{sec:met-models} and~\ref{sec:met-prompts} justify 
the model and prompt selection. Section~\ref{sec:met-scoring} details the 
log-probability scoring pipeline. Section~\ref{sec:met-metrics} defines the 
calibration and validity metrics. Section~\ref{sec:met-stats} presents the 
statistical analysis plan, including the recalibration and conformal 
prediction procedures.

\begin{figure}[htbp]
    \centering
    % [TODO: create a schematic diagram showing:
    %   Dataset (1015 pairs) 
    %      → 6 models × 2 prompts = 12 configurations
    %      → Log-probability scoring
    %      → Metrics (ECE, Brier, C-index)
    %      → Statistical tests (H1-H4) + Recalibration + Conformal]
    \caption{Overview of the experimental design. Six open-source LLMs are 
        evaluated with two prompt variants on 1{,}015 patient--criterion pairs, 
        producing twelve configurations tested against four pre-registered 
        hypotheses.}
    \label{fig:met-overview}
\end{figure}

% ----------------------------------------------------------------------------
\section{Research Question and Pre-registered Hypotheses}
\label{sec:met-hypotheses}

The main research question addressed in this work is stated as follows:

\begin{quote}
    \itshape
    Are the eligibility scores produced by large language models in 
    patient--clinical trial matching probabilistically reliable, and does 
    reinforcement learning from human feedback affect this reliability?
\end{quote}

This question is decomposed into four falsifiable hypotheses, each associated 
with a numerical rejection threshold defined prior to the experimental phase. 
Pre-registration of hypotheses is a standard safeguard against confirmation 
bias in empirical work~\cite{collins2024tripod} and provides a clear criterion 
for evaluating each finding.

Hypothesis~H1 tests whether current LLM pipelines are systematically 
miscalibrated. The rejection threshold is derived from Van~Calster 
et~al.~\cite{vancalster2019calibration}, who characterize an Expected 
Calibration Error (ECE) below $0.05$ as evidence of a well-calibrated 
predictive system in clinical contexts.

\begin{hypothesis}[H1: Poor calibration of current LLM pipelines]
    The Expected Calibration Error of an LLM-based eligibility scoring 
    pipeline exceeds the well-calibrated threshold of $0.05$:
    \[
        \ece(\text{pipeline}) > 0.05.
    \]
\end{hypothesis}

Hypothesis~H2 extends the finding of Kadavath et~al.~\cite{kadavath2022know}, 
who observed that RLHF post-training degrades calibration on general 
question-answering tasks, to the specific domain of clinical trial matching. 
It is the central original contribution of this thesis. The comparison uses 
a permutation test on the difference in ECE between paired Base and Instruct 
versions of the same model family, with Bonferroni correction across families 
and prompt variants.

\begin{hypothesis}[H2: RLHF degrades calibration]
    Within a given model family, the RLHF-aligned Instruct variant exhibits 
    worse calibration than its Base counterpart:
    \[
        \ece(\text{Instruct}) > \ece(\text{Base}).
    \]
\end{hypothesis}

Hypothesis~H3 tests whether, independently of absolute calibration, the 
model's scores respect the ordinal structure of the eligibility labels. 
Ordinal validity is measured by the concordance index (C-index), and the 
threshold of $0.70$ corresponds to conventional discrimination criteria used 
in clinical prediction modeling~\cite{vancalster2019calibration}.

\begin{hypothesis}[H3: Ordinal validity]
    The C-index between the LLM ordinal score and the expert ordinal label 
    exceeds the standard discrimination threshold:
    \[
        \operatorname{C-index} > 0.70.
    \]
\end{hypothesis}

Hypothesis~H4 addresses the practical value of confidence-based abstention. 
If confidence scores are meaningful, restricting predictions to the most 
confident subset should yield higher classification quality. The threshold 
of $+0.02$ on macro-F1 gain is chosen to reflect a clinically noticeable 
improvement without being trivially achievable.

\begin{hypothesis}[H4: Selective classification improves quality]
    Restricting predictions to the top 80\% most confident cases increases 
    macro-F1 by at least $0.02$:
    \[
        F_1(\text{coverage}=80\%) > F_1(\text{coverage}=100\%) + 0.02.
    \]
\end{hypothesis}

Together, these four hypotheses probe distinct aspects of probabilistic 
reliability: absolute calibration (H1), the effect of RLHF (H2), ordinal 
consistency (H3), and practical usability under abstention (H4). All 
rejection criteria and statistical tests were fixed prior to running the 
experiments and are reported without modification in Chapter~\ref{chap:results}.

% ----------------------------------------------------------------------------
\section{Data Preparation}
\label{sec:met-data}

The experimental evaluation is conducted on the TrialGPT-Criterion-Annotations 
dataset~\cite{jin2024trialgpt}, released alongside the TrialGPT system by 
Jin et al. To our knowledge, it is the only publicly available benchmark that 
provides multi-level expert annotations at the level of individual eligibility 
criteria, with an explicit distinction between inclusion and exclusion 
criteria. These properties are essential for a fine-grained calibration 
analysis: aggregating annotations at the trial level, as in TREC Clinical 
Trials or SIGIR benchmarks, would preclude the per-criterion probability 
estimation required by the present protocol.

\subsection{Dataset Composition}
\label{subsec:met-data-composition}

The dataset comprises 1{,}015 patient--criterion pairs. Each pair associates 
a synthetic patient note with a single eligibility criterion drawn from a 
real clinical trial. Two independent annotations are provided per pair: a 
consensus label from three physicians, and a label produced by GPT-4. Both 
annotations use the same six-level ordinal scale, described in 
Table~\ref{tab:met-labels}.

\begin{table}[htbp]
    \centering
    \caption{The six eligibility labels used in TrialGPT-Criterion-Annotations,
        organized by criterion type.}
    \label{tab:met-labels}
    \small
    \begin{tabular}{@{}p{2.5cm}p{2.5cm}p{7cm}@{}}
        \toprule
        \textbf{Label} & \textbf{Criterion type} & \textbf{Semantics} \\
        \midrule
        included & inclusion & The patient satisfies the inclusion criterion. \\
        not included & inclusion & The patient does not satisfy the criterion. \\
        excluded & exclusion & The patient has the excluding condition. \\
        not excluded & exclusion & The patient does not have the excluding condition (favorable). \\
        not enough information & both & The note does not permit a decision. \\
        not applicable & both & The criterion does not apply to this patient. \\
        \bottomrule
    \end{tabular}
\end{table}

Two of the six labels function as evasion or uncertainty markers rather than 
decisive clinical judgments (``not enough information'' and 
``not applicable''). We deliberately keep them in the label space to preserve 
the semantics of the original benchmark, while treating their behavior as an 
object of study in its own right (see Sections~\ref{sec:res-evasion} 
and~\ref{sec:res-mode-collapse}).

\subsection{Label Distribution}
\label{subsec:met-data-distribution}

The distribution of expert annotations is markedly imbalanced 
(Table~\ref{tab:met-label-freq}). ``Not excluded'' dominates at 47.7\%, 
reflecting the fact that most patients do not present the specific medical 
conditions targeted by exclusion criteria. ``Not enough information'' 
accounts for a further 28.9\%. The remaining 23.4\% of pairs, distributed 
across four labels, represent the cases requiring a truly discriminative 
clinical judgment. This asymmetry is critical for the interpretation of 
results in Chapter~\ref{chap:results}: aggregate metrics computed over the 
full 1{,}015 pairs can be dominated by the two majority classes and may 
mask complete failure on the discriminative subset.

\begin{table}[htbp]
    \centering
    \caption{Expert-label distribution in TrialGPT-Criterion-Annotations.}
    \label{tab:met-label-freq}
    \begin{tabular}{@{}lcc@{}}
        \toprule
        \textbf{Label} & \textbf{Count} & \textbf{Proportion} \\
        \midrule
        not excluded & 484 & 47.7\% \\
        not enough information & 293 & 28.9\% \\
        included & 74 & 7.3\% \\
        excluded & 55 & 5.4\% \\
        not applicable & 55 & 5.4\% \\
        not included & 54 & 5.3\% \\
        \midrule
        Total & 1{,}015 & 100.0\% \\
        \bottomrule
    \end{tabular}
\end{table}

\subsection{Textual Characteristics}
\label{subsec:met-data-text}

Patient notes are short and structured: their length ranges from 232 to 
1{,}447 characters with a median of 579 characters. This compactness is 
important for two reasons. First, it means that no aggressive truncation is 
required to fit the notes within LLM context windows, avoiding an additional 
source of variability. Second, short notes place a heavier reasoning burden 
on the model: information is often implicit, and the model must infer 
unstated conditions from the surrounding clinical context. Criterion texts 
are also short, with a median length of 58 characters, ranging from 
declarative statements (\emph{Age 18 to 75 years}) to more elaborate clauses 
involving numeric thresholds or drug interactions.

\subsection{Split for Recalibration Analysis}
\label{subsec:met-data-split}

The full dataset is used for the primary calibration audit (Chapter~\ref{chap:results}, 
Section~\ref{sec:res-h1}) since the models are queried without any prior 
fitting on this data. For the post-hoc recalibration analysis 
(Section~\ref{sec:res-recalibration}), we adopt a 70/30 random split 
stratified by criterion type. The 70\% subset is used to fit the isotonic 
regression, the 30\% subset to evaluate the recalibrated ECE. Stability of 
the recalibration gain is further verified through five-fold cross-validation 
using the same split ratio.

\subsection{Ethical and Reproducibility Considerations}
\label{subsec:met-data-ethics}

The dataset contains no real patient data: all notes are synthetic, 
constructed by the original authors to match the semantics of real 
recruitment scenarios without disclosing personal health information. No 
institutional review board approval was therefore required for the present 
work. The dataset is publicly available under a CC BY 4.0 license, and all 
prediction files produced in this study are archived in the project 
repository (see Appendix~\ref{app:code}) to support independent 
reproduction of the analyses.

% ----------------------------------------------------------------------------
% ----------------------------------------------------------------------------
\section{Model Selection}
\label{sec:met-models}

The choice of models under evaluation is the central design decision of this 
work. Rather than surveying a broad catalogue of open-source LLMs, we 
prioritize a small selection organized as a controlled experimental matrix. 
This design enables clean isolation of the effects under study --- 
particularly the effect of RLHF post-training --- while remaining feasible on 
the limited computational resources available.

\subsection{Selection Criteria}
\label{subsec:met-criteria}

Three criteria guide the selection of models. First, the model must be 
available in both a base and an instruction-tuned variant with a common 
pretraining, since testing hypothesis~H2 requires paired comparisons within 
the same model family. Second, the selection must span the two-by-two matrix 
of interest: generalist versus medically specialized pretraining, crossed 
with the presence or absence of RLHF post-training. Third, the model must 
fit into 16~GB of GPU memory when quantized to 4~bits, which is the memory 
budget of a free-tier T4 GPU on Kaggle. This constraint effectively caps the 
parameter count at 7--9 billion.

These three criteria eliminate several candidates that would otherwise be 
attractive. Models larger than 9~B parameters (Meditron-70B, Med42-v2-70B, 
Llama-3.3-70B) exceed the memory budget. Model families without a clean 
base--instruct pairing (Med42-v2 releases only the instruction-tuned 
variant) preclude the H2 comparison. Model families released without a 
medical counterpart (Gemma, DeepSeek) do not add information for our 
two-by-two matrix.

\subsection{The Six Selected Models}
\label{subsec:met-selected}

The six selected models, together with their scientific role in the 
experimental matrix, are summarized in Table~\ref{tab:met-models}.

\begin{table}[htbp]
    \centering
    \caption{Selected models and their scientific role in the two-by-two 
        experimental matrix.}
    \label{tab:met-models}
    \small
    \begin{tabular}{@{}p{4.5cm}p{4cm}p{4.5cm}@{}}
        \toprule
        \textbf{Model} & \textbf{Type} & \textbf{Scientific Role} \\
        \midrule
        \model{Mistral-7B-v0.3} & Generalist, no RLHF & H2 control \\
        \model{Mistral-7B-Instruct-v0.3} & Generalist, with RLHF & H2 test \\
        \model{Llama-3.1-8B} & Generalist, no RLHF & H2 replication \\
        \model{Llama-3.1-8B-Instruct} & Generalist, with RLHF & H2 replication \\
        \model{Meditron-7B} & Medical, no RLHF & Pure specialization effect \\
        \model{BioMistral-7B} & Medical, inherited RLHF & Validation vs.\ Labrak et al.\ (2024) \\
        \bottomrule
    \end{tabular}
\end{table}

\model{Mistral-7B-v0.3} and its instruction-tuned counterpart 
\model{Mistral-7B-Instruct-v0.3} form the primary base--instruct pair for 
testing~H2. The two models share the same pretraining corpus and 
architecture; the Instruct variant differs only through subsequent 
supervised fine-tuning and RLHF. This clean isolation of the RLHF stage 
makes them the reference comparison for our central hypothesis.

\model{Llama-3.1-8B} and \model{Llama-3.1-8B-Instruct} play the same role 
for a second, independent model family. Testing H2 on two families rather 
than one rules out the possibility that any observed effect is specific to 
a single training pipeline. Llama-3.1 also brings a different tokenizer, a 
larger context window, and a different alignment methodology, which 
strengthens the generality of the conclusions.

\model{Meditron-7B}~\cite{chen2023meditron} is a foundation model 
pretrained on 46{,}000 clinical practice guidelines and a large corpus of 
PubMed abstracts, starting from Llama-2 weights. Crucially, no official 
instruction-tuned variant of Meditron-7B has been released; the base model 
is the only publicly available version. Meditron therefore isolates the 
effect of medical pretraining without any confound with RLHF, providing a 
direct test of whether medical specialization improves calibration in the 
absence of instruction alignment.

\model{BioMistral-7B}~\cite{labrak2024biomistral} occupies the opposite 
corner of the two-by-two matrix. It is derived from 
Mistral-7B-Instruct-v0.1 through continued pretraining on PubMed Central 
Open Access texts. It therefore inherits the RLHF alignment of its 
Mistral-Instruct parent while adding a medical specialization stage. 
BioMistral is also the only open-source medical LLM with published ECE 
values, which provides an external validation point for our measurement 
methodology.

Together, these six models populate the two-by-two matrix completely 
(Table~\ref{tab:met-matrix}). Each of the three scientific questions 
associated with the matrix --- the effect of RLHF, the effect of medical 
specialization, and their interaction --- can be tested through paired 
comparisons within this design.

\begin{table}[htbp]
    \centering
    \caption{The two-by-two experimental matrix of selected models.}
    \label{tab:met-matrix}
    \begin{tabular}{@{}lcc@{}}
        \toprule
         & \textbf{Without RLHF} & \textbf{With RLHF} \\
        \midrule
        \textbf{Generalist} & 
            \model{Mistral-7B} & 
            \model{Mistral-7B-Instruct} \\
         & 
            \model{Llama-3.1-8B} & 
            \model{Llama-3.1-8B-Instruct} \\
        \midrule
        \textbf{Medical} & 
            \model{Meditron-7B} & 
            \model{BioMistral-7B} \\
        \bottomrule
    \end{tabular}
\end{table}

\subsection{Excluded Models}
\label{subsec:met-excluded}

Several candidate models were considered and deliberately excluded from the 
final selection. \model{Qwen2.5-7B} and \model{DeepSeek-V2-Lite} were 
excluded as redundant additions to the generalist axis: they would not 
provide new evidence for H2 beyond what Mistral and Llama already establish. 
\model{Gemma-2-9B} was excluded for the same reason, and additionally 
because its license restricts commercial redistribution of derivative 
model outputs, complicating reproducibility.

Among medical LLMs, \model{Med42-v2}~\cite{christophe2024med42} and 
\model{OpenBioLLM-8B} were excluded because their base versions are not 
publicly released, precluding a within-family H2 comparison. 
\model{HuatuoGPT-o1}~\cite{chen2024huatuogpt} was excluded because its 
training corpus is primarily Chinese, which introduces a confound with the 
English-only nature of the TrialGPT-Criterion benchmark. Larger medical 
models (\model{Meditron-70B}, \model{Med42-v2-70B}, \model{Aloe-70B}) all 
exceed the memory budget of a free-tier T4 GPU, even at 4-bit quantization.

The rationale for this restrictive selection is that a small number of 
well-matched models tested on the same protocol yields more interpretable 
conclusions than a broader benchmark that would necessarily mix confounding 
factors. Extending the analysis to additional families is a natural 
direction for future work, discussed in Chapter~\ref{chap:conclusion}.
% ----------------------------------------------------------------------------
% ----------------------------------------------------------------------------
\section{Prompt Design}
\label{sec:met-prompts}

The prompt used to elicit predictions from a large language model is not a 
neutral component of an evaluation pipeline: it is a design choice that 
directly affects the model's outputs. Any calibration audit that fixes a 
single prompt inevitably conflates two effects --- the properties of the 
model and the properties of the prompt --- which cannot be separated post 
hoc. To address this ambiguity, we evaluate each of the six models under 
two distinct prompt variants, treating the prompt as an experimental variable 
rather than a fixed convention.

\subsection{Rationale for Two Prompt Variants}
\label{subsec:met-two-prompts}

The two prompts are designed to represent the two extremes of a spectrum of 
elicitation strategies commonly encountered in practice. The first 
(Prompt~A) is intentionally minimal: it presents the classification task and 
the label vocabulary without any explanation, in the manner of a naive user 
querying the model. The second (Prompt~B) is enriched with definitions, 
examples, and explicit reasoning rules, in the manner of an expert operator 
who invests time in prompt engineering. By evaluating each model with both 
prompts, we obtain four independent measurement points per model family for 
hypothesis~H2 (Base--Instruct comparison at Prompt~A, Base--Instruct 
comparison at Prompt~B), doubling the statistical power of the RLHF test.

This design also enables three additional analyses beyond the pre-registered 
hypotheses. First, the effect of prompt design itself becomes measurable 
within each model. Second, the interaction between RLHF and prompt design 
becomes testable: some models may be more sensitive to prompt enrichment 
than others. Third, the design places our results in a stronger position 
with respect to external validity: a finding replicated under two different 
prompts is more robust than a finding measured under a single one.

\subsection{Prompt A: Generic Baseline}
\label{subsec:met-prompt-a}

Prompt~A is a minimalist prompt containing only the classification 
instruction, the input note, the criterion text, and the list of candidate 
labels. It does not define the labels, does not provide examples, and does 
not include reasoning rules. Its purpose is to establish a baseline that 
corresponds to how a first-time user would query the model, and to provide 
a reference point against which enriched prompts can be compared.

The complete text of Prompt~A is given in Appendix~\ref{app:prompts}. Its 
essential structure is:

\begin{quote}
\itshape
Classify whether this patient meets the eligibility criterion.\\
Patient: \textrm{[note]}\\
Criterion (\textrm{[type]}): \textrm{[criterion]}\\
Choose exactly one label from: included, not included, excluded, 
not excluded, not enough information, not applicable.\\
Answer:
\end{quote}

Prompt~A is applied identically to all six models. For instruction-tuned 
models, the prompt is wrapped in the model's native chat template. For base 
models, the prompt is presented as a plain-text continuation task without 
special tokens.

\subsection{Prompt B: Adapted with Expert Guidance}
\label{subsec:met-prompt-b}

Prompt~B is an adapted prompt containing definitions of the six labels, 
one worked example per label, and a critical inference rule for exclusion 
criteria that addresses a subtle failure mode observed in preliminary 
experiments. The full text is reproduced in Appendix~\ref{app:prompts}.

Three design choices distinguish Prompt~B from Prompt~A. First, each of the 
six labels is defined with an explicit semantic description. For example, 
``not included'' is characterized as ``The patient note contains 
information showing the patient does NOT meet this inclusion criterion,'' 
which distinguishes it from ``not enough information'' (where the note is 
silent). Second, a one-sentence example is provided for each label, drawn 
from realistic clinical situations. Third, a specific inference rule is 
stated for exclusion criteria:

\begin{quote}
\itshape
If the criterion names a specific disease, condition, or medical history, 
and the note does not mention it, the default is ``not excluded'' (the 
patient most likely does not have it). Reserve ``not enough information'' 
only for criteria requiring a specific lab value or measurement that is 
truly missing.
\end{quote}

This rule directly targets the tendency of language models to produce 
``not enough information'' for any exclusion criterion whose target 
condition is not explicitly discussed in the patient note --- a pattern 
that mechanically converts the absence of a mention into an expression of 
epistemic uncertainty, when clinically it should be treated as evidence of 
absence.

\subsection{Adaptations for Base vs.\ Instruct Models}
\label{subsec:met-base-instruct}

Base models and instruction-tuned models expect prompts in different 
formats, and applying the same prompt verbatim to both would create an 
artificial handicap for the base variant. To ensure a fair comparison, 
we adapt the surface form of each prompt to the target model class while 
preserving its semantic content.

For instruction-tuned models, prompts are formatted using the model's 
native chat template, which delimits user and assistant roles with 
model-specific special tokens (\texttt{[INST]}~for Mistral, 
\texttt{<|start\_header\_id|>}~for Llama-3). This is the format for which 
these models were aligned and where they behave as expected.

For base models, no chat template is applied. Instead, the prompt includes 
one or more few-shot examples that establish the desired input-output 
pattern by demonstration. Prompt~A for base models includes a single 
example; Prompt~B includes three examples covering distinct label types. 
This asymmetry reflects a genuine limitation of base models --- they do 
not natively follow instructions --- and treating it as an evaluation 
artifact would misrepresent how these models are actually used in 
practice. The introduction of few-shot examples is nevertheless a 
potential confounder for cross-family comparisons, and we discuss it 
explicitly in Chapter~\ref{chap:conclusion}.

The complete text of all four prompt variants (Prompt~A / Prompt~B, 
Instruct / Base) is provided in Appendix~\ref{app:prompts} to support 
independent reproduction of the pipeline.

% ----------------------------------------------------------------------------
% ----------------------------------------------------------------------------
\section{Scoring Pipeline: Log-Probability Extraction}
\label{sec:met-scoring}

Standard practice for querying a large language model on a classification 
task is to sample a textual output and post-process it to recover a label. 
This approach has two limitations for a calibration audit. First, the 
sampled output is a point estimate: it discards the underlying probability 
distribution over candidate labels, which is precisely the object of 
interest for calibration analysis. Second, generation is stochastic, 
introducing an additional source of variability that would complicate the 
reproducibility of downstream statistics. We therefore adopt a 
log-probability scoring approach that directly reads out the model's 
probability distribution over the candidate labels, without any generation 
step.

\subsection{Principle of Log-Probability Scoring}
\label{subsec:met-scoring-principle}

For a given patient--criterion pair, we first construct the prompt as 
described in Section~\ref{sec:met-prompts}. Rather than sampling a 
continuation, we consider each candidate label as a hypothetical completion 
of the prompt. For each label $\ell$, we compute the mean log-probability 
of its token sequence conditional on the prompt:
\begin{equation}
    s_\ell = \frac{1}{|\ell|} \sum_{t=1}^{|\ell|} \log p(\ell_t \mid \text{prompt}, \ell_{<t})
    \label{eq:label-logprob}
\end{equation}
where $\ell_t$ is the $t$-th token of the label and $|\ell|$ is the total 
number of tokens. Length normalization by $|\ell|$ is necessary because 
labels have unequal token lengths in the tokenizer: without it, longer 
labels such as ``not enough information'' would be systematically 
penalized against shorter ones such as ``included''.

The resulting scores $\{s_\ell\}$ are then normalized through a softmax to 
obtain a proper probability distribution $\pi$ over the candidate labels:
\begin{equation}
    \pi_\ell = \frac{\exp(s_\ell)}{\sum_{\ell' \in \mathcal{L}} \exp(s_{\ell'})}
    \label{eq:softmax}
\end{equation}
The predicted label $\hat{y}$ is the argmax of $\pi$, and the associated 
confidence $\hat{p}$ is the maximum probability. Crucially, the full 
distribution $\pi$ is retained for downstream analyses, including 
recalibration, conformal prediction, and the simulation of constrained 
label spaces (Section~\ref{sec:res-mode-collapse}).

\subsection{Restriction of the Candidate Set by Criterion Type}
\label{subsec:met-scoring-candidates}

The six labels described in Section~\ref{sec:met-data} split into two 
groups by criterion type. For inclusion criteria, the semantically valid 
labels are \emph{included}, \emph{not included}, \emph{not enough 
information}, and \emph{not applicable}. For exclusion criteria, they are 
\emph{excluded}, \emph{not excluded}, \emph{not enough information}, and 
\emph{not applicable}. In both cases, four labels are candidates and two 
would be semantically nonsensical (e.g., predicting ``included'' for an 
exclusion criterion). We therefore restrict the softmax in 
Equation~\ref{eq:softmax} to the four candidates appropriate to the 
criterion type at hand. This restriction preserves the ordinal structure 
of the task and prevents the model from being penalized for probability 
mass placed on labels that could not have been correct in the first place.

\subsection{Algorithmic Summary}
\label{subsec:met-scoring-algorithm}

Algorithm~\ref{alg:scoring} summarizes the full scoring procedure applied 
to each patient--criterion pair.

\begin{algorithm}[htbp]
\caption{Log-probability scoring for eligibility label prediction}
\label{alg:scoring}
\begin{algorithmic}[1]
    \Require Patient note $n$, criterion $c$, criterion type $t$
    \Ensure Predicted label $\hat{y}$, confidence $\hat{p}$, distribution $\pi$
    \State Build prompt from $n$, $c$, $t$
    \State $\mathcal{L} \gets$ candidate labels according to $t$
    \ForAll{$\ell \in \mathcal{L}$}
        \State $s_\ell \gets $ mean log-probability of tokens of $\ell$
    \EndFor
    \State $\pi \gets \operatorname{softmax}(s)$
    \State $\hat{y} \gets \arg\max_{\ell} \pi_\ell$
    \State $\hat{p} \gets \pi_{\hat{y}}$
    \State \Return $\hat{y}, \hat{p}, \pi$
\end{algorithmic}
\end{algorithm}

Because no sampling is involved, the procedure is fully deterministic: 
running the same model on the same input twice yields identical scores. 
This determinism is important for the reproducibility of the entire 
downstream analysis, including the bootstrap confidence intervals and 
permutation tests described in Section~\ref{sec:met-stats}.

\subsection{Implementation Notes}
\label{subsec:met-scoring-implementation}

The scoring loop is implemented using the standard causal language modeling 
interface of the HuggingFace \texttt{transformers} library. For a given 
prompt, we tokenize the concatenation of the prompt and each candidate 
label separately, run a single forward pass per concatenation, and read 
out the token log-probabilities from the model's logits after applying a 
log-softmax over the vocabulary. This approach requires four forward 
passes per patient--criterion pair (one per candidate label), which is 
approximately twice as expensive as free-form generation of a short 
answer. In practice, on a T4 GPU with 4-bit quantization, the four passes 
complete in approximately three seconds per pair for a 7-billion-parameter 
model, and the full evaluation of 1{,}015 pairs takes fifty to seventy 
minutes depending on the model.

% ----------------------------------------------------------------------------
% ----------------------------------------------------------------------------
\section{Quantization and Infrastructure}
\label{sec:met-infrastructure}

Running 7-to-9-billion-parameter models on a free-tier T4 GPU (16~GB of 
memory) requires reducing the memory footprint of the model weights. We 
apply 4-bit NF4 quantization~\cite{dettmers2023qlora} through the 
\texttt{bitsandbytes} library, which represents each weight as a 4-bit 
value drawn from a distribution matched to the empirical distribution of 
neural network weights. This reduces the memory required to store a 
7-billion-parameter model from approximately 14~GB (in half-precision) to 
approximately 4~GB, leaving sufficient headroom for the key-value cache 
during inference.

Prior work has shown that 4-bit NF4 quantization introduces negligible 
degradation on downstream evaluation benchmarks compared to full 
precision~\cite{dettmers2023qlora}, including on tasks that involve 
long-form generation or fine-grained probability estimation. We rely on 
this evidence to justify the use of NF4 quantization for our calibration 
audit, while acknowledging that a small residual effect of quantization 
on absolute ECE values cannot be entirely ruled out. Since Hypothesis~H2 
and the associated Bonferroni-corrected permutation tests compare 
\emph{differences} in ECE between paired variants evaluated under the 
same quantization regime, any systematic bias introduced by quantization 
cancels out and does not affect the conclusions.

All experiments were run on Kaggle notebooks with a single Tesla~T4 GPU. 
The compute quota (30~GPU-hours per week on the free tier) was sufficient 
to complete the twelve runs and all downstream analyses within the 
experimental phase of the project.

% ----------------------------------------------------------------------------
% ----------------------------------------------------------------------------
\section{Evaluation Metrics}
\label{sec:met-metrics}

The evaluation combines calibration metrics, classification metrics, 
ordinal validity metrics, and behavioral metrics. Each captures a distinct 
aspect of probabilistic reliability, and no single metric is sufficient on 
its own. This section defines every metric used in Chapter~\ref{chap:results}.

\subsection{Expected Calibration Error (ECE)}
\label{subsec:met-ece}

The Expected Calibration Error, introduced by Naeini et~al.\ and 
popularized by Guo et~al.~\cite{guo2017calibration}, measures the average 
absolute gap between predicted confidence and observed accuracy across 
confidence bins:

\begin{equation}
    \ece = \sum_{k=1}^{K} \frac{|B_k|}{n}
    \left| \operatorname{acc}(B_k) - \operatorname{conf}(B_k) \right|
    \label{eq:ece}
\end{equation}

where predictions are partitioned into $K$ equal-width bins $B_1, \ldots, 
B_K$ based on their confidence, $n$ is the total number of predictions, 
$\operatorname{acc}(B_k)$ is the empirical accuracy of predictions in 
bin $B_k$, and $\operatorname{conf}(B_k)$ is the average confidence of 
predictions in the same bin. A perfectly calibrated model satisfies 
$\operatorname{acc}(B_k) = \operatorname{conf}(B_k)$ for all bins, yielding 
an ECE of zero. We use $K = 10$ equal-width bins on $[0, 1]$.

\subsection{Brier Score}
\label{subsec:met-brier}

The Brier score~\cite{brier1950} measures the mean squared error between 
predicted probability and the true binary outcome (correct or incorrect):

\begin{equation}
    \operatorname{Brier} = \frac{1}{n} \sum_{i=1}^{n}
    \left( p_i - y_i \right)^2
    \label{eq:brier}
\end{equation}

where $p_i$ is the predicted confidence and $y_i \in \{0, 1\}$ indicates 
whether the prediction was correct. Unlike ECE, which measures calibration 
alone, the Brier score jointly reflects calibration and sharpness: a model 
that is both accurate and confident achieves a lower Brier score than one 
that is only well calibrated but uninformative.

\subsection{Classification Metrics}
\label{subsec:met-classification}

We report accuracy as the proportion of predictions matching the expert 
label, and macro-F1 as the unweighted mean of per-class F1 scores. Macro-F1 
is critical in the presence of the severe class imbalance documented in 
Section~\ref{subsec:met-data-distribution}: a model that always predicts 
the majority class achieves an accuracy of 47.7\% but a macro-F1 close to 
zero on the minority classes. This distinction becomes central in 
Chapter~\ref{chap:results} when comparing model behavior across cases of 
varying clinical difficulty.

\subsection{C-index for Ordinal Validity}
\label{subsec:met-cindex}

Because the eligibility labels form an ordinal scale, we complement 
calibration and classification metrics with the concordance index 
(C-index), which measures the probability that, for a randomly chosen 
pair of cases with different expert labels, the model correctly ranks 
them:

\begin{equation}
    \operatorname{C-index} = \frac{
        |\{(i, j) : y_i < y_j \text{ and } s_i < s_j\}|
    }{
        |\{(i, j) : y_i \neq y_j\}|
    }
    \label{eq:cindex}
\end{equation}

where $y_i, y_j$ are ordinal expert labels of cases $i$ and $j$, and 
$s_i, s_j$ are the model's scalar scores obtained by mapping the 
predicted label and confidence onto the ordinal scale. A C-index of 1 
indicates perfect ranking, 0.5 indicates random ranking, and the 
conventional threshold of 0.70 is used in clinical prediction modeling as 
evidence of acceptable discrimination~\cite{vancalster2019calibration}.

\subsection{Weighted Kappa for Expert-Model Agreement}
\label{subsec:met-kappa}

To measure the agreement between physician annotations and GPT-4 
annotations in the pre-LLM audit (Section~\ref{sec:res-audit}), we use 
Cohen's weighted kappa with quadratic weights~\cite{cohen1968}. Weighted 
kappa extends the standard kappa coefficient by penalizing disagreements 
proportionally to their ordinal distance: a confusion between adjacent 
labels (e.g., ``included'' and ``not enough information'') counts less 
than a confusion between distant labels (e.g., ``included'' and 
``excluded''). We report values in the range $[-1, 1]$, with 1 
indicating perfect agreement, 0 indicating chance agreement, and 
negative values indicating systematic disagreement.

\subsection{Overconfidence and Prediction Diversity}
\label{subsec:met-behavioral}

Beyond calibration and classification, two behavioral metrics quantify 
how a model uses its probability distribution.

The mean overconfidence is defined as the difference between the average 
confidence and the average accuracy:

\begin{equation}
    \operatorname{overconfidence} = \bar{p} - \operatorname{acc}
    \label{eq:overconf}
\end{equation}

where $\bar{p}$ is the mean predicted confidence across the evaluation 
set. A positive value indicates that the model tends to overstate its 
certainty, a negative value that it understates it. Unlike ECE, which is 
a magnitude, overconfidence is a signed quantity that reveals the 
direction of the miscalibration.

Prediction diversity is measured as the normalized Shannon entropy of 
the empirical label distribution:

\begin{equation}
    \operatorname{diversity} = 
    \frac{
        -\sum_{\ell \in \mathcal{L}} f_\ell \log_2 f_\ell
    }{
        \log_2 |\mathcal{L}|
    }
    \label{eq:diversity}
\end{equation}

where $f_\ell$ is the empirical frequency of label $\ell$ in the model's 
predictions and $\mathcal{L}$ is the label vocabulary. A diversity of 0 
indicates that the model predicts a single label; a diversity of 1 
indicates a uniform distribution over all labels. This metric is central 
to detecting the mode collapse behavior documented in 
Section~\ref{sec:res-mode-collapse}, which the calibration and 
classification metrics can mask.

\subsection{Composite Utility Score}
\label{subsec:met-composite}

To provide a single rank-ordering of the twelve configurations that 
reflects clinical usability, we propose a composite utility score 
combining calibration quality, discrimination, and prediction diversity:

\begin{equation}
    \operatorname{utility} = F_1^{\text{macro}} \times (1 - \ece) \times 
    \operatorname{diversity}
    \label{eq:composite}
\end{equation}

The three factors are multiplicative rather than additive so that a 
failure on any single dimension zeroes out the score. A model with 
excellent calibration but complete mode collapse (diversity $= 0$) 
receives a utility score of zero, correctly capturing its lack of 
practical value. This composite score complements rather than replaces 
the individual metrics, which remain the primary object of analysis.

\subsection{Reliability Diagrams}
\label{subsec:met-reliab}

Reliability diagrams complement ECE by visualizing the calibration curve 
directly. For each of the $K$ bins used in the ECE computation, we plot 
the average confidence on the horizontal axis and the observed accuracy 
on the vertical axis. A perfectly calibrated model produces points that 
lie on the diagonal $y = x$. Points below the diagonal indicate 
overconfidence; points above indicate underconfidence.

\subsection{Bootstrap Confidence Intervals}
\label{subsec:met-bootstrap}

For every point estimate of ECE, Brier score, and C-index, we report a 
95\% confidence interval obtained by 1{,}000 bootstrap resamples of the 
prediction set. Bootstrap confidence intervals do not assume any specific 
distributional form and are appropriate for the non-Gaussian statistics 
under study. They also provide a natural way to assess whether a given 
model's ECE genuinely exceeds the 0.05 threshold of Hypothesis~H1, or 
whether the excess is within sampling noise.

% ----------------------------------------------------------------------------
\section{Statistical Analysis Plan}
\label{sec:met-stats}

The four pre-registered hypotheses described in 
Section~\ref{sec:met-hypotheses} are tested according to a fixed analysis 
plan. This section describes the tests, the correction for multiple 
comparisons, and the additional analyses (post-hoc recalibration and 
conformal prediction) that complement the primary hypotheses.

\subsection{Pre-registration of Hypotheses}
\label{subsec:met-preregistration}

All four rejection criteria (H1: $\ece > 0.05$; H2: 
$\ece(\text{Instruct}) > \ece(\text{Base})$; H3: 
$\operatorname{C-index} > 0.70$; H4: $F_1(80\%) > F_1(100\%) + 0.02$) 
were fixed before the twelve model runs were executed. This 
pre-registration protects the analysis against post-hoc adjustment of 
the criteria to fit the observed results and follows the recommendations 
of the TRIPOD+AI reporting guidelines~\cite{collins2024tripod}. Results 
are reported without modification of the rejection criteria, and any 
findings not covered by these four hypotheses are reported separately as 
original observations (see Chapter~\ref{chap:results}).

\subsection{Permutation Tests}
\label{subsec:met-permutation}

Hypothesis~H2 compares the ECE of paired Base and Instruct variants of 
the same model family under the same prompt. Rather than assume any 
particular distributional form for the ECE difference, we use a 
permutation test: for 1{,}000 random permutations of the model labels 
(Base vs.\ Instruct) within each pair, we recompute the ECE difference 
and count the fraction of permutations in which the observed difference 
is equalled or exceeded. This fraction is the one-sided p-value for the 
null hypothesis that the two variants have the same underlying 
calibration. The one-sided formulation is appropriate because 
Hypothesis~H2 is directional: we test whether Instruct is worse than 
Base, not merely different from it.

\subsection{Bonferroni Correction}
\label{subsec:met-bonferroni}

Hypothesis~H2 is tested four times (Mistral--Prompt~A, Mistral--Prompt~B, 
Llama--Prompt~A, Llama--Prompt~B) and further comparisons for medical 
specialization contribute additional tests. To control the family-wise 
error rate across the eight planned comparisons, we apply a Bonferroni 
correction with $\alpha_{\text{corrected}} = 0.05 / 8 = 0.00625$. This 
correction is conservative but appropriate given the small number of 
comparisons and the desire for a stringent test of the central 
hypothesis.

\subsection{Post-hoc Recalibration Methods}
\label{subsec:met-recalibration}

Beyond the primary calibration audit, we evaluate three post-hoc 
recalibration methods that could correct miscalibration without retraining 
the model. Temperature scaling~\cite{guo2017calibration} applies a single 
learned scalar $T$ to the logits before the softmax, effectively 
softening (or sharpening) the confidence distribution. Isotonic 
regression fits a monotone piecewise-constant function mapping 
uncorrected confidences to corrected ones; being non-parametric, it can 
correct arbitrary shapes of miscalibration but requires more data than 
temperature scaling. Platt scaling fits a logistic function of the 
uncorrected confidence, offering a compromise between the flexibility of 
isotonic regression and the parsimony of temperature scaling.

Each method is fitted on the 70\% training subset and evaluated on the 
30\% held-out subset (see Section~\ref{subsec:met-data-split}). To assess 
whether the recalibration gains are stable across different data splits, 
we additionally perform a 5-fold cross-validation, reporting the mean 
and standard deviation of the ECE reduction over the five folds.

\subsection{Conformal Prediction}
\label{subsec:met-conformal}

To complement point-estimate calibration, we apply split conformal 
prediction~\cite{vovk2005}, a distribution-free framework that produces 
prediction sets with a guaranteed marginal coverage rate. Given a 
target coverage level $1 - \alpha$ (we use $\alpha = 0.10$, i.e., 90\% 
coverage), the framework calibrates a threshold on the non-conformity 
scores of a held-out calibration set and returns, for each test case, 
the set of labels whose non-conformity is below that threshold. The 
size of the prediction set is an informative quantity: large sets 
indicate high aleatoric uncertainty, while small sets indicate high 
model confidence. We report both empirical coverage (which should match 
the target 90\%) and average set size.

% ----------------------------------------------------------------------------
\section{Reproducibility}
\label{sec:met-reproducibility}

All code, prompts, and prediction files produced in this study are 
archived in the project repository listed in Appendix~\ref{app:code}. The 
code is organized as a single Kaggle notebook that reproduces the twelve 
model runs from scratch, together with a set of Python scripts that 
perform the downstream analyses (calibration audit, hypothesis tests, 
recalibration, and conformal prediction) from the archived prediction 
files. Random seeds are fixed throughout the pipeline: the 4-bit 
quantization is deterministic, the log-probability scoring is 
deterministic by construction (Section~\ref{sec:met-scoring}), and the 
bootstrap and permutation tests use a fixed seed to ensure that 
statistical results are exactly reproducible.

The complete list of hardware and software versions used in the 
experiments, together with the exact commit hashes of the notebook and 
the analysis scripts, is provided in Appendix~\ref{app:code}. Any 
independent researcher with access to a T4 GPU (or equivalent) can 
reproduce the twelve model runs in approximately ten to fifteen hours of 
GPU time, and the downstream analyses in a few additional minutes on 
CPU.